\section{Finite strategies for bounded quantifier alternation}
\label{sec:strategy}

A finite theorem often has a bounded interactive form rather than a single
existential certificate.  For example,
\[
  \forall x<B_0(q)\;\exists y<B_1(q,x)\;
  \forall z<B_2(q,x,y)\;R(q,x,y,z).
\]
It is incorrect to replace the existential moves by an infinite choice
function.  For a fixed bounded formula, however, the complete finite strategy
is itself one certificate.

\subsection{Strategy tables}

For each existential node, enumerate the preceding universal history in mixed
radix and store one bounded response in a table.  A strategy code \(s\) then
has primitive-recursive accessors \(\mathsf{move}_j(s,h)\).  The verifier checks:
\begin{enumerate}[label=(\roman*)]
\item every move is below its history-dependent bound;
\item for every complete universal history, the matrix bit is true after
      inserting the stored moves.
\end{enumerate}
The second condition is a bounded fold over all histories, represented by a
zero failure sum or least-counterexample search.

\begin{theorem}[Strategy soundness, fixed shape]
\label{thm:strategy-sound}
For every fixed alternation shape and reflected quantifier-free matrix, there
is a specialized primitive-recursive verifier
\(\mathsf{verify}_\alpha(q,s)\) and an iPRA derivation
\[
  \mathsf{verify}_\alpha(q,s)=1
   \;\longrightarrow\;
  \Phi_\alpha(q).
\]
\end{theorem}

\begin{proof}
Proceed by structural induction on the fixed alternation shape at the Lean
metalevel.  An existential node reads the stored move, proves its bound from
the verifier, and introduces the existential.  A universal node opens an
arbitrary bounded challenge and specializes the verifier's no-failure result
to its history code.  At a leaf apply matrix reflection.  The bounded-history
specialization is obtained from the least-counterexample or failure-sum
reflection of \cref{sec:reflection}; no object induction contains a quantified
invariant.
\end{proof}

\subsection{Completeness for finite decidable games}

When all move bounds are explicit and the leaf matrix is decidable, backward
induction computes a winning bit and a least winning response at each
existential node.  Because the alternation shape is fixed, this is a
primitive-recursive program family.  The object proof uses nested bounded
search and quantifier-free reflection.

\begin{theorem}[Finite strategy completeness, schema]
For each fixed bounded alternation shape \(\alpha\), iPRA proves that if the
bounded formula \(\Phi_\alpha(q)\) holds, then an accepted finite strategy code
exists.  Equivalently, the formula is interderivable with the existence of an
accepted strategy certificate.
\end{theorem}

\begin{proof}[Proof architecture]
Induct metatheoretically on \(\alpha\).  At an existential node, retain the
formula's witness in local scope and combine it with the recursively obtained
sub-strategy; if a named table is desired, use the explicit move and table
bounds to select the least accepted combined code.  At a universal node,
construct a table of sub-strategies for all finitely many challenges.  This
table construction is a primitive recursion on the challenge bound whose
invariant states, in reflected form, that all completed entries verify.  Since
\(\alpha\) is fixed, all table accessors and verifiers are specialized syntax.
A complete Lean proof should first implement the \(\forall\exists\) case and
then abstract the recurring fold.
\end{proof}

The completeness schema is a formalization target rather than a claimed
kernel-checked declaration in this archive.  Its human proof uses only the
bounded search and trace lemmas already isolated above.

\subsection{Dual certificates and counterstrategies}

For diagnostic use, compile a dual certificate for failure: a universal
challenge or counterstrategy exposing a rejected leaf.  Since the finite game
is decidable, the winning and losing bits are complementary.  A packet can
therefore separate three concerns:
\begin{itemize}
\item a small checker for a positive strategy;
\item a small checker for a negative counterexample;
\item the producer theorem deciding which certificate exists.
\end{itemize}
This is useful for finite combinatorics, covering arguments, and bounded
algebraic searches.  It does not imply determinacy for infinite games or an
unbounded choice principle.
